% ============================================================================
% ARTÍCULO 1 — RCTA (Revista Colombiana de Tecnologías de Avanzada)
% Clasificación ordinal del desgaste de brocas CNC mediante features acústicas
% Autor: Oscar Iván Ayala Ortiz
% ============================================================================
\documentclass[10pt,twocolumn]{article}

% ── Paquetes ─────────────────────────────────────────────────────────────────
\usepackage[utf8]{inputenc}
\usepackage[T1]{fontenc}
\usepackage[spanish,es-tabla]{babel}
\usepackage{mathptmx}               % Times New Roman
\usepackage[top=2.5cm,bottom=2.5cm,left=2cm,right=2cm]{geometry}
\usepackage{graphicx}
\usepackage{booktabs}
\usepackage{amsmath,amssymb}
\usepackage{multirow}
\usepackage{caption}
\usepackage{subcaption}
\usepackage{hyperref}
\usepackage{xcolor}
\usepackage{float}
\usepackage{array}
\usepackage{tabularx}
\usepackage{siunitx}
\sisetup{output-decimal-marker={,}}

\hypersetup{colorlinks=true,linkcolor=blue,citecolor=blue,urlcolor=blue}

% ── Ruta de figuras ──────────────────────────────────────────────────────────
\graphicspath{{../results/article1_figures/}}

% ── Metadata ─────────────────────────────────────────────────────────────────
\title{\textbf{Clasificación ordinal del desgaste de brocas en taladrado CNC
mediante análisis acústico: sensibilidad al umbral de etiquetado}}
% ----------------- Fondo (Plantilla) en la capa de fondo -----------------
\usepackage{tikz}
\usepackage{eso-pic}
\newcommand{\CIETABackground}{Plantilla_CIETA.pdf} % RENOMBRAR archivo a Plantilla_CIETA.pdf
\AddToShipoutPictureBG{
  \AtPageLowerLeft{
    \includegraphics[width=\paperwidth,height=\paperheight]{\CIETABackground}
  }
}


\author{
  Jorge Enrique Meneses Flórez$^{1}$,\,
  Nicolás Orejarena Osorio$^{2}$,\,
  Oscar Iván Ayala Ortiz$^{3}$\\[4pt]
  {\small Escuela de Ingeniería Mecánica, Universidad Industrial de Santander, La Universidad, Bucaramanga, Colombia$^{1,2,3}$}\\[4pt]
  {\small \texttt{E-mail:} \texttt{jmeneses@uis.edu.co}$^{1}$, \texttt{norejarena@correo.uts.edu.co}$^{2}$, \texttt{ayalaortizoscarivan@gmail.com}$^{3}$}
}

\date{}

% ============================================================================
\begin{document}
\maketitle

% ── RESUMEN ──────────────────────────────────────────────────────────────────
\begin{abstract}
Se presenta un sistema de clasificación ordinal del desgaste de brocas
helicoidales en taladrado CNC de acero AISI~4140, orientado al
mantenimiento predictivo y basado exclusivamente en señales acústicas
capturadas con micrófonos de bajo costo. Se emplea el método de
descomposición ordinal de Frank y Hall con máquinas de vectores de
soporte (SVM) y bosques aleatorios (RF), utilizando 10 \textit{features}
acústicas seleccionadas estadísticamente. La contribución central es un
análisis de sensibilidad al umbral de etiquetado, barriendo el porcentaje
de vida útil que define la clase ``desgastado'' desde el 60\% hasta el
97\%. Los resultados muestran que las \textit{features} acústicas
\textit{hand-crafted} detectan eficazmente el fallo catastrófico
(umbral~$\geq$\,97\%: $F_1$~macro = 0{,}46, exactitud adyacente = 98{,}6\%),
pero su capacidad disminuye para desgaste progresivo temprano
(umbral~$\geq$\,60\%: $F_1$~macro = 0{,}30, exactitud adyacente = 83{,}4\%).
Un hallazgo clave es que la exactitud adyacente permanece superior al 83\%
en todo el rango de umbrales, indicando que los errores de clasificación
nunca exceden un paso ordinal. Mediante una etapa de refinamiento que
incorpora ingeniería de \textit{features} derivadas, optimización de
hiperparámetros (kernel lineal, $C = 5$, pesos de clase asimétricos) y
calibración de umbrales de decisión, la exactitud adyacente de la SVM al
umbral del 75\% alcanza 98{,}5\% y la exactitud exacta 60{,}2\%,
superando a la configuración base (85{,}9\% y 50{,}3\%,
respectivamente). El análisis de generalización cruzada entre siete
experimentos independientes revela que el tipo de micrófono es la
principal fuente de \textit{domain shift}. La reducción de ruido espectral
(\textit{spectral gating}) mejora la exactitud adyacente en test de
84{,}7\% a 90{,}1\% y reduce el \textit{domain shift} entre tipos de
micrófono ($F_1$: 0{,}18 $\rightarrow$ 0{,}33; exactitud: 23\%
$\rightarrow$ 59\% en validación con micrófono dinámico). Estos
resultados delimitan el alcance operacional del monitoreo acústico pasivo
con \textit{features} \textit{hand-crafted}, con implicaciones directas
para la implementación en sistemas de monitoreo en línea de bajo costo
computacional, y motivan el uso de representaciones espectrales aprendidas
para la detección de desgaste gradual.

\vspace{6pt}
\noindent\textbf{Palabras clave:} desgaste de herramientas, taladrado CNC,
clasificación ordinal, análisis acústico, SVM, sensibilidad al umbral,
mantenimiento predictivo, Industria~4.0.
\end{abstract}

% ── ABSTRACT ──────────────────────────────────────────────────────────────────
\renewcommand{\abstractname}{Abstract}
\begin{abstract}
An ordinal classification system for drill bit wear in CNC drilling of
AISI~4140 steel is presented, aimed at predictive maintenance and based
exclusively on acoustic signals captured with low-cost microphones. The
Frank and Hall ordinal decomposition method is employed with support
vector machines (SVM) and random forests (RF), using 10 statistically
selected acoustic features. The central contribution is a labeling
threshold sensitivity analysis, sweeping the tool life percentage that
defines the ``worn'' class from 60\% to 97\%. Results show that
hand-crafted acoustic features effectively detect catastrophic failure
(threshold~$\geq$\,97\%: macro~$F_1$ = 0.46, adjacent accuracy = 98.6\%),
but their discriminative power decreases for early progressive wear
(threshold~$\geq$\,60\%: macro~$F_1$ = 0.30, adjacent accuracy = 83.4\%).
A key finding is that adjacent accuracy remains above 83\% across all
thresholds, indicating that classification errors never exceed one ordinal
step. A refinement stage incorporating derived feature engineering,
hyperparameter optimization (linear kernel, $C = 5$, asymmetric class
weights), and decision threshold calibration raises adjacent accuracy at
the 75\% threshold to 98.5\% and exact accuracy to 60.2\%, surpassing
the baseline configuration (85.9\% and 50.3\%, respectively).
Cross-experiment generalization analysis across seven independent
experiments reveals that microphone type is the primary source of domain
shift. Spectral gating noise reduction improves test adjacent accuracy
from 84.7\% to 90.1\% and reduces cross-microphone domain shift
(validation $F_1$: 0.18 $\rightarrow$ 0.33; accuracy: 23\%
$\rightarrow$ 59\%).

\vspace{6pt}
\noindent\textbf{Keywords:} tool wear, CNC drilling, ordinal classification,
acoustic analysis, SVM, threshold sensitivity, predictive maintenance,
Industry~4.0.
\end{abstract}

% ============================================================================
\section{Introducción}
\label{sec:intro}

El monitoreo de la condición de herramientas de corte es un componente
esencial en la manufactura moderna, donde la detección oportuna del desgaste
permite reducir costos de producción, mejorar la calidad dimensional de las
piezas y prevenir fallas catastróficas que comprometan la integridad de la
máquina y la seguridad del operario~\cite{groover2007}. En operaciones de
taladrado CNC, el desgaste progresivo del filo de corte modifica las fuerzas
de corte, las vibraciones mecánicas y la emisión acústica del
proceso~\cite{bhuiyan2016}.

El enfoque acústico para el monitoreo de herramientas ha ganado atención por
ser no invasivo, de bajo costo y fácil de integrar en entornos industriales
existentes~\cite{charoenprasit2018, orejarena2014, meneses2020}. Si bien
los enfoques basados en aprendizaje profundo (\textit{deep learning}) han
demostrado alta precisión en tareas de clasificación de señales, su costo
computacional limita su despliegue en sistemas embebidos de borde
(\textit{edge computing}) propios de la manufactura
inteligente~\cite{groover2007}. En contraste, los modelos basados en
\textit{features} acústicas \textit{hand-crafted} con clasificadores
ligeros (SVM, RF) permiten inferencia en tiempo real con recursos
computacionales mínimos, lo que los hace particularmente adecuados para
aplicaciones de mantenimiento predictivo en el contexto de la
Industria~4.0.

Sin embargo, la mayoría de los trabajos previos abordan el problema como
clasificación binaria (nueva vs.\ desgastada) o categórica sin considerar
la naturaleza progresiva del desgaste~\cite{banda2024}. Esta simplificación
ignora información ordinal valiosa: confundir una herramienta nueva con una
moderadamente desgastada tiene consecuencias operacionales distintas a
confundirla con una herramienta en fallo.

Un segundo problema frecuente es la definición del umbral que separa las
clases de desgaste. La mayoría de los estudios utilizan un umbral fijo
(típicamente basado en el criterio de flanco VB según ISO~8688 o
ISO~3685~\cite{iso8688}) sin explorar la sensibilidad de los resultados a
esta elección. En escenarios donde no se dispone de medición directa de VB
--- como cuando se utilizan exclusivamente señales acústicas --- el umbral
se define sobre el porcentaje de vida útil consumida, lo que introduce una
decisión metodológica cuyo impacto rara vez se cuantifica.

El presente trabajo aborda ambas limitaciones. Se formula el problema como
clasificación ordinal, es decir, una tarea en la que las clases tienen un orden 
natural entre sí y, por tanto, no se consideran categorías totalmente independientes.
 En este caso se clasifican en tres clases (sin desgaste, medianamente 
 desgastado,desgastado) utilizando el método de Frank y Hall~\cite{frank2001} con SVM y
RF. La contribución central es la identificacion del porcentaje de vida util
correspondiente a la zona desgaste. Para ello se iteraron diferetes rangos de etiquetado
hasta encontrar la precisión ajustada, la cual ajusta los problemas de desbalance debido a
que no se tiene la misma cantidad de datos en cada estado. De esta manera se obtuvo que la
definición de desgastado corresponde de un 60\% hasta el 98.5\% de la vida útil. 
Este barrido permite mapear los
límites operacionales del monitoreo acústico pasivo con \textit{features}
\textit{hand-crafted} y establecer recomendaciones sobre qué umbrales son
viables para el estado medianamente desgastado versus detección del estado desgastado.


Las contribuciones específicas son: (1)~formulación ordinal del desgaste de
brocas con métricas apropiadas (MAE ordinal, exactitud adyacente);
(2)~análisis de sensibilidad al umbral con curva completa de 9 umbrales;
(3)~evaluación de generalización cruzada entre 7 experimentos independientes;
(4)~estudio de ablación de \textit{features} que revela que un subconjunto
de 7 \textit{features} supera al conjunto completo de 10;
(5)~evaluación del impacto de la reducción de ruido espectral como
preprocesamiento para mitigar el \textit{domain shift} entre tipos de
micrófono.

% ============================================================================
\section{Materiales y métodos}
\label{sec:metodos}

% ── 2.1 Setup experimental ──────────────────────────────────────────────────
\subsection{Configuración experimental}

Los ensayos de taladrado se realizaron en un centro de mecanizado CNC
Leadwell V-20 (recorridos: $510 \times 350 \times \SI{400}{mm}$,
velocidad máxima del husillo: \SI{8000}{rpm}). El material de trabajo fue
acero AISI~4140 bonificado (28--32~HRC), en discos de $6'' \times 2''$
sujetos con copa periférica. Se utilizaron brocas helicoidales Dormer A100
de \SI{6}{mm} y \SI{8}{mm} de diámetro, con trayectoria en espiral o
lineal según el ensayo. Todos los ensayos emplearon refrigerante
(taladrina) y agujeros ciegos de \SI{25}{mm} de profundidad.

Se realizaron siete ensayos independientes (E1--E7), cuyas condiciones se
resumen en la Tabla~\ref{tab:experimentos}. Se usaron 3 diferentes trayectorias,
lineal, expiral y mixto. Los ensayos E1--E4 se
condujeron hasta falla completa de la broca (tres clases de desgaste),
mientras que E5--E7 registraron desgaste progresivo sin alcanzar falla
(dos clases).

\begin{table}[htbp]
\centering
\caption{Condiciones experimentales de los siete ensayos.}
\label{tab:experimentos}
\footnotesize
\begin{tabular}{@{}lccccr@{}}
\toprule
Exp. & $\varnothing$ (mm) & G-code & Micrófono & Tipo & $n$ \\
\midrule
E1 & 8 & mixto    & dinámico     & Hasta falla & 84 \\
E2 & 8 & espiral  & dinámico     & Hasta falla & 128 \\
E3 & 6 & lineal   & condensador  & Hasta falla & 583 \\
E4 & 6 & espiral  & condensador  & Hasta falla & 1194 \\
E5 & 8 & lineal   & condensador  & Progresivo  & 234 \\
E6 & 8 & lineal   & condensador  & Progresivo  & 234 \\
E7 & 6 & espiral  & condensador  & Progresivo  & 148 \\
\midrule
\multicolumn{5}{l}{Total} & 2565(definir n de la tabla) \\

\bottomrule
\end{tabular}
\end{table}

Los parámetros de corte se calcularon según las recomendaciones del
fabricante: velocidad de corte $V_c = \SI{20}{m/min}$, avance
$f = \SI{0.109}{mm/rev}$ (6~mm) y $f = \SI{0.138}{mm/rev}$ (8~mm),
resultando en velocidades del husillo de \SI{1061}{rpm} y \SI{796}{rpm},
respectivamente.

% ── 2.2 Sistema de adquisición ──────────────────────────────────────────────
\subsection{Sistema de adquisición de datos}

La adquisición se realizó con un chasis NI~cDAQ-9174 equipado con un
módulo NI-9234 (4 canales, \SI{51.2}{kS/s}, 24~bits, filtro
\textit{anti-aliasing} integrado). Se emplearon dos micrófonos en
configuraciones alternadas según el ensayo~\cite{orejarena2014}:

\begin{itemize}
  \item \textbf{M1}: TXL~235 TOPP~AUDIO (dinámico), utilizado en E1 y E2.
  \item \textbf{M2}: PG81 SHURE (condensador, cardioide), utilizado en
    E3--E7.
\end{itemize}

\noindent Esta asignación implica que el tipo de micrófono constituye una
covariable confundida con el diseño experimental, aspecto que se analiza
en la Sección~\ref{sec:discusion}.

Las señales se adquirieron a una frecuencia de muestreo de
$f_s = \SI{51.2}{kS/s}$ (ancho de banda útil hasta \SI{25.6}{kHz}),
en formato \textit{Technical Data Management Streaming} (TDMS), y se convirtieron a WAV PCM de 16~bits sin compresión
para preservar la integridad de los transitorios
acústicos~\cite{macphail2024}.

% ── 2.3 Reducción de ruido ────────────────────────────────────────────────
\subsection{Preprocesamiento: reducción de ruido espectral}
\label{sec:noisereduce}

Para evaluar el impacto del ruido de fondo en la clasificación, se
implementó un \textit{pipeline} de programación para la reducción del ruido 
basado en el método de filtrado
\textit{spectral gating} estacionario. Se construyó un perfil de ruido
combinó grabaciones de referencia, tales como
ruido del husillo sin taladrado ($\approx$\SI{180}{min} disponibles),
husillo con refrigerante ($\approx$\SI{46}{min}) y ambiente del
laboratorio ($\approx$\SI{48}{min}), todas adquiridas con el mismo sistema
NI-9234 a \SI{51.2}{kS/s} segmentadas cada 120s para efectos de entrenamiento.

El filtrado se aplicó a los 2635 archivos originales WAV mediante la librería
\texttt{noisereduce}~\cite{sainburg2020} con los parámetros: modo
estacionario (\texttt{stationary=True}), proporción de atenuación
$\alpha = 0{,}85$, $n_{\text{FFT}} = 2048$ y la señal de referencia como
perfil de ruido anteriormente combinado. La señal filtrada conserva sus niveles relativos
(sin normalización de pico) para no alterar las \textit{features} (características)
dependientes de energía. De estos se reextrajeron las mismas 7 \textit{features} (características) con mas influencia en los 
resultados de predicción a la hora de entrenar los modelos empleando un estudio de ablación.

% ── 2.4 Extracción de features ──────────────────────────────────────────────
\subsection{Extracción de \textit{features} (características) acústicas del modelo que predice}

Se extrajeron 21 \textit{features} ancla que cubren tres dominios
físicos: timbre (centroide espectral, MFCC, planitud, contraste, entropía,
chroma, ratio armónico/percusivo), energía (RMS, factor de cresta, energía
Mel y wavelet) y temporal (duración, tasa de \textit{onsets}, ZCR). A
partir de un proceso de selección multi-criterio --- bootstrap de
estabilidad en top-$N$, ANOVA/Kruskal-Wallis, análisis de no-redundancia
e información mutua --- se seleccionaron las 10 \textit{features} de la
Tabla~\ref{tab:features}.

\begin{table}[htbp]
\centering
\caption{\textit{Features} acústicas seleccionadas (línea base).}
\label{tab:features}
\footnotesize
\begin{tabular}{@{}rlll@{}}
\toprule
\# & Feature & Dominio & Interpretación física \\
\midrule
1 & spectral\_contrast & Timbre & Diferencia picos/valles espectrales \\
2 & crest\_factor      & Energía & Ratio pico/RMS (impulsividad) \\
3 & chroma\_std        & Timbre & Variabilidad cromática \\
4 & zcr                & Temporal & Cruces por cero (ruido de banda ancha) \\
5 & spectral\_entropy  & Timbre & Desorden espectral \\
6 & centroid\_mean     & Timbre & Centro de masa espectral \\
7 & harmonic\_perc.    & Timbre & Ratio armónico/percusivo \\
8 & onset\_rate        & Temporal & Frecuencia de transitorios \\
9 & spectral\_flatness & Timbre & Tonalidad vs.\ ruido \\
10 & duration\_s       & Temporal & Duración del ciclo \\
\bottomrule
\end{tabular}
\end{table}

% ── 2.4 Estrategia de etiquetado ────────────────────────────────────────────
\subsection{Estrategia de etiquetado con barrido de umbrales}
\label{sec:etiquetado}

Al no disponer de mediciones directas de desgaste de flanco (VB), las
etiquetas se asignan según el porcentaje de vida útil consumida, definido
como la razón entre el número de agujero actual y el número total de
agujeros realizados con cada broca. Se definen tres clases ordinales:

\begin{equation}
y_i = \begin{cases}
  \text{sin desgaste}        & \text{si } p_i \leq \tau_{\text{sin}} \\
  \text{med.\ desgastado}    & \text{si } \tau_{\text{sin}} < p_i < \tau_{\text{des}} \\
  \text{desgastado}           & \text{si } p_i \geq \tau_{\text{des}}
\end{cases}
\label{eq:etiquetado}
\end{equation}

\noindent donde $p_i = h_i / h_{\max}$ es la fracción de vida útil consumida
para la muestra $i$, $\tau_{\text{sin}} = 0{,}15$ (fijo) y
$\tau_{\text{des}}$ es el umbral que se varía. Los ensayos E5--E7
(progresivos, sin falla) conservan sus etiquetas originales, asignadas
mediante análisis de tendencia de \textit{features} acústicas (RMS,
centroide, ZCR) a lo largo del ensayo, siguiendo el enfoque correlacional
de Meneses y Peña~\cite{meneses2020}.

La contribución central de este trabajo es barrer $\tau_{\text{des}}$ en
nueve valores: 60\%, 65\%, 70\%, 75\%, 80\%, 85\%, 90\%,
95\% y 97\%. Para cada umbral se reetiqueta el dataset completo, se
reentrenan los modelos y se evalúa en el conjunto de test, permitiendo
construir una curva de sensibilidad que mapea la capacidad discriminativa
de las \textit{features} acústicas en función del grado de desgaste.

% ── 2.5 Clasificación ordinal ───────────────────────────────────────────────
\subsection{Clasificación ordinal: método de Frank y Hall}

Se emplea el método de Frank y Hall~\cite{frank2001}, que descompone un
problema ordinal de $K$ clases en $K-1$ clasificadores binarios. Para
tres clases ($k \in \{0, 1, 2\}$, donde 0 = sin desgaste, 1 = medianamente
desgastado, 2 = desgastado), se entrenan dos clasificadores:

\begin{align}
C_1&: P(y \geq 1) \quad \text{``¿Hay algún desgaste?''} \\
C_2&: P(y \geq 2) \quad \text{``¿El desgaste es severo?''}
\end{align}

Las probabilidades se combinan con restricción de monotonicidad
($\hat{P}_2 \leq \hat{P}_1$) para obtener la distribución ordinal:

\begin{align}
P(\text{sin}) &= 1 - \hat{P}_1 \\
P(\text{med}) &= \hat{P}_1 - \hat{P}_2 \\
P(\text{des}) &= \hat{P}_2
\end{align}

\noindent y la clase predicha es $\hat{y} = \arg\max_k P(k)$.

Cada clasificador binario $C_j$ es una SVM con kernel RBF, optimizada
mediante \textit{GridSearchCV} ($C \in \{1, 10\}$,
$\gamma \in \{\text{scale}, 0{,}1\}$) con validación cruzada
\textit{StratifiedGroupKFold} de 3 particiones, agrupando por experimento
para evitar fuga de información entre ensayos del mismo grupo
experimental. Esta estrategia garantiza que las muestras del mismo
experimento nunca aparezcan simultáneamente en entrenamiento y validación
interna~\cite{montgomery2013}. Se utiliza
\texttt{class\_weight='balanced'} para compensar el desbalance de clases.
Como comparación, se entrena un RF ordinal con la misma descomposición de
Frank y Hall (200 árboles, \texttt{class\_weight='balanced'}).

Adicionalmente, se realizó una etapa de refinamiento que incluyó: (i)~7
\textit{features} derivadas por combinación de las originales (ratios
espectrales, indicadores de estabilidad tonal y energía), (ii)~búsqueda
de hiperparámetros extendida sobre 30 configuraciones (kernels lineal y
RBF, $C \in \{0{,}5;\; 1;\; 5;\; 10;\; 50\}$, selección de $k \in
\{10;\; 15;\; 20;\; \text{todas}\}$ \textit{features} por información
mutua, pesos de clase asimétricos), y (iii)~calibración de los umbrales
de decisión $t_1$ y $t_2$ del clasificador ordinal (barrido en
$[0{,}15;\; 0{,}70]$ con paso de $0{,}05$), optimizando un \textit{score}
compuesto: $0{,}4 \times \text{adj.\ acc.} + 0{,}3 \times \text{bal.\
acc.} + 0{,}3 \times \text{exact.\ acc.}$

% ── 2.6 Diseño experimental ─────────────────────────────────────────────────
\subsection{Diseño experimental y métricas}
\label{sec:diseno}

La partición del dataset respeta la independencia experimental:

\begin{itemize}
  \item \textbf{Test}: E3 completo (583 muestras, 6~mm, condensador).
  \item \textbf{Validación}: E2 completo (128 muestras, 8~mm, dinámico).
  \item \textbf{Entrenamiento}: E1 + E4 + E5 + E6 + E7 (2224 muestras,
    incluyendo augmentados acústicos).
\end{itemize}

Se seleccionó E3 como test por ser el ensayo más extenso con falla
completa (583 muestras, tres clases), y E2 como validación por presentar
el mayor \textit{domain shift} (micrófono dinámico, 8~mm), permitiendo
evaluar robustez. La augmentación acústica del entrenamiento incluyó
\textit{time stretching}, adición de ruido gaussiano y desplazamiento de
\textit{pitch}, aplicadas exclusivamente a muestras de la clase
``desgastado'' para mitigar el desbalance (90 muestras augmentadas).
Ningún dato augmentado aparece en validación ni test, y los augmentados
derivados de E3 se excluyen del entrenamiento para evitar fuga de
información.

Las métricas reportadas incluyen:
\begin{itemize}
  \item \textbf{Macro $F_1$}: promedio no ponderado del $F_1$ por clase.
  \item \textbf{Exactitud adyacente}: fracción de predicciones con error
    $\leq 1$ paso ordinal (ec.~\ref{eq:adj_acc}).
  \item \textbf{MAE ordinal}: error absoluto medio entre clases predichas
    y verdaderas (ec.~\ref{eq:mae}).
  \item \textbf{Exactitud exacta}: fracción de predicciones correctas.
\end{itemize}

\begin{align}
\text{Adj.\ Acc.} &= \frac{1}{n} \sum_{i=1}^{n} \mathbb{1}[|\hat{y}_i - y_i| \leq 1]
\label{eq:adj_acc} \\
\text{MAE}_{\text{ord}} &= \frac{1}{n} \sum_{i=1}^{n} |\hat{y}_i - y_i|
\label{eq:mae}
\end{align}

Los intervalos de confianza al 95\% se estiman mediante
\textit{bootstrap} con 1000 remuestreos del conjunto de test.

% ============================================================================
\section{Resultados}
\label{sec:resultados}

% ── 3.1 Análisis exploratorio ───────────────────────────────────────────────
\subsection{Análisis exploratorio}

La Fig.~\ref{fig:pca} muestra la proyección PCA bidimensional del
conjunto de test (E3) sobre las 10 \textit{features} seleccionadas. Los
dos primeros componentes principales explican el 67{,}9\% de la
varianza (PC1: 41{,}8\%, PC2: 26{,}1\%). Se observa que las muestras
de clase ``sin desgaste'' se separan parcialmente del grupo principal,
mientras que ``medianamente desgastado'' y ``desgastado'' presentan
solapamiento considerable, consistente con la frontera difusa entre
desgaste progresivo.

\begin{figure}[htbp]
  \centering
  \includegraphics[width=\columnwidth]{pca_2d.pdf}
  \caption{Proyección PCA bidimensional del test set E3 (10
    \textit{features}, reetiquetado al 75\%).}
  \label{fig:pca}
\end{figure}

La distribución de clases por experimento (Fig.~\ref{fig:class_dist})
evidencia que E4 domina el dataset (1194 muestras) y que los ensayos
progresivos (E5--E7) carecen de clase ``desgastado'', lo que limita su
contribución al entrenamiento del clasificador $C_2$.

\begin{figure}[htbp]
  \centering
  \includegraphics[width=\columnwidth]{class_distribution.pdf}
  \caption{Distribución de clases por experimento (reetiquetado
    $\tau_{\text{des}} = 75\%$).}
  \label{fig:class_dist}
\end{figure}

El ranking de información mutua (Fig.~\ref{fig:mi}) para el conjunto
completo de 26 \textit{features} muestra que las \textit{features} de
variabilidad espectral (\texttt{spectral\_bandwidth\_std},
\texttt{rolloff\_std}) ganan relevancia con umbrales bajos (desgaste
gradual), mientras que \texttt{spectral\_contrast\_mean} domina a umbral
alto (97\%, fallo catastrófico). Entre las 10 \textit{features} del
baseline, \texttt{spectral\_contrast\_mean} mantiene la mayor información
mutua con la etiqueta de desgaste en todos los umbrales, seguida de
\texttt{centroid\_mean} y \texttt{zcr}.

\begin{figure}[htbp]
  \centering
  \includegraphics[width=\columnwidth]{mi_ranking.pdf}
  \caption{Ranking de \textit{features} por información mutua
    (reetiquetado al 75\%).}
  \label{fig:mi}
\end{figure}

% ── 3.2 Sensibilidad al umbral ─────────────────────────────────────────────
\subsection{Curva de sensibilidad al umbral}
\label{sec:threshold}

La Fig.~\ref{fig:threshold} y la Tabla~\ref{tab:threshold} presentan el
resultado principal del trabajo: la curva de sensibilidad al umbral de
etiquetado.

\begin{figure*}[htbp]
  \centering
  \includegraphics[width=\textwidth]{threshold_sensitivity_curve.pdf}
  \caption{Sensibilidad al umbral de etiquetado. (a)~Macro $F_1$ vs.\
    umbral. (b)~Exactitud adyacente vs.\ umbral. (c)~$F_1$ por clase
    (SVM). (d)~Número de muestras de desgaste en test y MAE ordinal.}
  \label{fig:threshold}
\end{figure*}

\begin{table*}[htbp]
\centering
\caption{Métricas de clasificación ordinal por umbral de etiquetado
  ($\tau_{\text{sin}} = 15\%$, test = E3).}
\label{tab:threshold}
\footnotesize
\begin{tabular}{@{}S[table-format=2.0]rrrrrrrrr@{}}
\toprule
{$\tau_{\text{des}}$ (\%)} &
{$n_{\text{des}}^{\text{test}}$} &
\multicolumn{4}{c}{SVM ordinal} &
\multicolumn{4}{c}{RF ordinal} \\
\cmidrule(lr){3-6} \cmidrule(lr){7-10}
& & {$F_1$} & {Acc.} & {Adj.} & {MAE} &
{$F_1$} & {Acc.} & {Adj.} & {MAE} \\
\midrule
60  & 230 & 0,304 & 0,381 & 0,834 & 0,786 & 0,369 & 0,437 & 0,909 & 0,654 \\
65  & 201 & 0,321 & 0,424 & 0,839 & 0,738 & 0,371 & 0,470 & 0,904 & 0,626 \\
70  & 172 & 0,348 & 0,475 & 0,839 & 0,686 & 0,398 & 0,533 & 0,926 & 0,540 \\
75  & 144 & 0,358 & 0,503 & 0,859 & 0,638 & 0,408 & 0,568 & 0,933 & 0,499 \\
80  & 115 & 0,363 & 0,520 & 0,887 & 0,593 & 0,419 & 0,581 & 0,942 & 0,477 \\
85  &  86 & 0,361 & 0,528 & 0,926 & 0,545 & 0,443 & 0,623 & 0,969 & 0,408 \\
90  &  58 & 0,386 & 0,551 & 0,954 & 0,496 & 0,443 & 0,638 & 0,978 & 0,384 \\
95  &  29 & 0,380 & 0,573 & 0,978 & 0,449 & 0,473 & 0,678 & 0,995 & 0,328 \\
97  &  18 & 0,463 & 0,583 & 0,986 & 0,431 & 0,437 & 0,705 & 0,995 & 0,300 \\
\bottomrule
\end{tabular}
\end{table*}

Se identifican tres hallazgos principales:

\textbf{1) El $F_1$ macro crece con el umbral.} La SVM pasa de $F_1 = 0{,}30$
(60\%) a $F_1 = 0{,}46$ (97\%), un incremento del 53\%. Esto confirma
que las \textit{features} acústicas \textit{hand-crafted} detectan más
fácilmente el fallo catastrófico (chirping, chatter) que el desgaste
gradual de flanco, cuya firma acústica es más sutil.

\textbf{2) La exactitud adyacente permanece alta.} Incluso al umbral más
exigente (60\%), la exactitud adyacente de la SVM es 83{,}4\% y la del
RF 90{,}9\%. A partir del 85\%, supera el 92\% (SVM) y 96\% (RF).
Esto significa que \textit{los errores de clasificación nunca saltan dos
clases}: un error típico confunde ``medianamente desgastado'' con
``desgastado'' o viceversa, pero nunca ``sin desgaste'' con ``desgastado''.

\textbf{3) El RF supera a la SVM base en casi todos los umbrales.} La ventaja
del RF es consistente en exactitud adyacente ($+4$ a $+8$ puntos
porcentuales) y $F_1$ macro ($+0{,}02$ a $+0{,}09$), excepto a 97\%
donde la SVM logra $F_1 = 0{,}46$ versus $F_1 = 0{,}44$ del RF. Esta
superioridad se explica porque el RF, al promediar múltiples árboles de
decisión con submuestreo aleatorio, es inherentemente más robusto al
desbalance de clases y a la variabilidad del espacio de \textit{features},
mientras que la SVM RBF con umbrales de decisión por defecto tiende a
concentrar sus predicciones en la clase mayoritaria. Sin embargo, tras la
etapa de refinamiento (Sección~\ref{sec:refinamiento}), la SVM optimizada
alcanza una exactitud adyacente de 98{,}5\% al umbral del 75\%,
superando al RF base (93{,}3\%) en este punto de operación.

Los intervalos de confianza \textit{bootstrap} (95\%, 1000 remuestreos)
para el modelo SVM entrenado con etiquetas originales (equivalente al
umbral 97\% con 4 muestras de desgaste en test) son:
$F_1 = 0{,}54\; [0{,}38;\; 0{,}75]$,
exactitud $= 94{,}0\%\; [92{,}1;\; 95{,}7]$,
exactitud adyacente $= 100\%$, MAE $= 0{,}06\; [0{,}04;\; 0{,}08]$.
El intervalo amplio del $F_1$ refleja el soporte estadístico limitado de la
clase ``desgastado'' a este umbral. Para el reetiquetado al 75\% (144
muestras de desgaste), el IC se estrecha: $F_1 = 0{,}35\; [0{,}32;\; 0{,}39]$.

% ── 3.3 Análisis por clase ──────────────────────────────────────────────────
\subsection{Análisis por clase}

El panel~(c) de la Fig.~\ref{fig:threshold} descompone el $F_1$ por clase
para la SVM. La clase ``medianamente desgastado'' mantiene $F_1 > 0{,}55$
en todos los umbrales gracias a su prevalencia en el test set. La clase
``sin desgaste'' es estable pero baja ($F_1 \approx 0{,}27\text{--}0{,}31$),
sugiriendo que el modelo confunde muestras iniciales con desgaste moderado.
La clase ``desgastado'' es la más sensible al umbral: su $F_1$ crece de
0{,}04 (60\%) a 0{,}37 (97\%), reflejando que solo el fallo catastrófico
produce una firma acústica claramente distinguible.

% ── 3.4 Ablación de features ───────────────────────────────────────────────
\subsection{Estudio de ablación}

La Fig.~\ref{fig:ablation} presenta el estudio de ablación con subconjuntos
incrementales de \textit{features} ordenadas por información mutua.

\begin{figure}[htbp]
  \centering
  \includegraphics[width=\columnwidth]{ablation_curve.pdf}
  \caption{Ablación de \textit{features}. (a)~$F_1$ vs.\ número de
    \textit{features} (top-$k$ por MI). (b)~Cambio en $F_1$ al remover
    cada \textit{feature} del top-10 (umbral 97\%).}
  \label{fig:ablation}
\end{figure}

Un resultado notable es que \textbf{el subconjunto top-7 supera al top-10}
en los tres umbrales evaluados (75\%: 0{,}367 vs.\ 0{,}358; 85\%:
0{,}378 vs.\ 0{,}353; 97\%: 0{,}441 vs.\ 0{,}406). Las tres
\textit{features} eliminadas (\texttt{onset\_rate},
\texttt{spectral\_flatness\_mean}, \texttt{duration\_s}) aportan ruido que
degrada la discriminación de la SVM en espacios de alta dimensionalidad
relativa.

El análisis \textit{leave-one-out} a 97\% (panel~b) confirma que
\texttt{crest\_factor} es perjudicial: su eliminación \textit{mejora} el
$F_1$ en $+0{,}064$. Similarmente, \texttt{centroid\_mean} mejora el $F_1$
en $+0{,}057$ al ser removido. Únicamente
\texttt{spectral\_contrast\_mean} muestra una contribución neta negativa
al ser eliminado ($-0{,}022$), confirmando su rol como la \textit{feature}
más discriminativa para desgaste severo.

% ── 3.5 Interpretabilidad: SHAP ─────────────────────────────────────────────
\subsection{Interpretabilidad: análisis SHAP}

La Fig.~\ref{fig:shap} presenta los diagramas \textit{beeswarm} de valores
SHAP~\cite{li2024} para los dos clasificadores ordinales ($C_1$ y $C_2$)
del modelo entrenado con umbral original (97\%), usando
\textit{KernelExplainer} con 200 muestras de fondo y 300 de evaluación.

\begin{figure}[htbp]
  \centering
  \begin{subfigure}[b]{0.48\columnwidth}
    \includegraphics[width=\linewidth]{shap_beeswarm_C1.pdf}
    \caption{$C_1$: P($y \geq 1$)}
  \end{subfigure}\hfill
  \begin{subfigure}[b]{0.48\columnwidth}
    \includegraphics[width=\linewidth]{shap_beeswarm_C2.pdf}
    \caption{$C_2$: P($y \geq 2$)}
  \end{subfigure}
  \caption{SHAP \textit{beeswarm} para los clasificadores ordinales.
    Los valores SHAP cuantifican el impacto de cada \textit{feature}
    en la predicción individual.}
  \label{fig:shap}
\end{figure}

Los valores SHAP revelan que \texttt{spectral\_contrast\_mean} domina la
discriminación en $C_2$ (P(\text{desgaste severo})), consistente con el
resultado de ablación. En $C_1$, \texttt{zcr} y \texttt{centroid\_mean}
tienen mayor peso relativo, reflejando que la presencia de desgaste
incipiente se manifiesta principalmente en el contenido espectral de
alta frecuencia. Los valores SHAP bajos (escala $\approx 0{,}05$) son
esperados dado que la SVM opera en espacio escalado y los
\textit{features} individuales tienen contribuciones distribuidas.

% ── 3.6 Impacto de la reducción de ruido ──────────────────────────────────
\subsection{Impacto de la reducción de ruido espectral}
\label{sec:noise_results}

La Tabla~\ref{tab:noise_comparison} compara el rendimiento del modelo SVM
ordinal (top-7 \textit{features}) entrenado sobre \textit{features}
originales versus \textit{features} extraídas de WAVs procesados con
\textit{spectral gating}.

\begin{table}[htbp]
\centering
\caption{Comparativa: SVM ordinal sin filtro vs.\ con reducción de
  ruido (\textit{spectral gating}, $\alpha=0{,}85$). Top-7 features.}
\label{tab:noise_comparison}
\footnotesize
\begin{tabular}{@{}llrrrr@{}}
\toprule
Fuente & Split & {$F_1$} & {Acc.} & {Adj.} & {MAE} \\
\midrule
Original  & val  & 0{,}180 & 0{,}234 & 0{,}914 & 0{,}852 \\
Filtrado  & val  & 0{,}331 & 0{,}586 & 1{,}000 & 0{,}414 \\
\midrule
Original  & test & 0{,}338 & 0{,}461 & 0{,}847 & 0{,}691 \\
Filtrado  & test & 0{,}342 & 0{,}489 & 0{,}901 & 0{,}611 \\
\bottomrule
\end{tabular}
\end{table}

En el conjunto de test (E3, condensador), el filtrado produce una mejora
modesta pero consistente: $+0{,}004$ en $F_1$ macro, $+5{,}3$ puntos en
exactitud adyacente (de 84{,}7\% a 90{,}1\%) y $-0{,}08$ en MAE ordinal.
El hallazgo más relevante se observa en validación (E2, micrófono
dinámico), donde el filtrado reduce drásticamente el \textit{domain shift}:
el $F_1$ macro pasa de 0{,}18 a 0{,}33 ($+83{,}9\%$ relativo), la
exactitud de 23{,}4\% a 58{,}6\% ($+150{,}9\%$ relativo) y la exactitud
adyacente alcanza 100\%. Esto indica que una fracción significativa del
\textit{domain shift} entre tipos de micrófono se debe a diferencias en
el piso de ruido captado por cada transductor, y que el \textit{spectral
gating} las atenúa al normalizar el contenido no informativo.

% ── 3.7 Refinamiento del modelo ──────────────────────────────────────────
\subsection{Refinamiento: ingeniería de \textit{features} y calibración}
\label{sec:refinamiento}

La etapa de refinamiento descrita en la Sección~\ref{sec:metodos} se
evaluó sobre el reetiquetado al 75\% ($\tau_{\text{des}} = 0{,}75$). La
Tabla~\ref{tab:refinamiento} compara el modelo base (SVM RBF,
\texttt{class\_weight='balanced'}, umbrales 0{,}50/0{,}50) con el modelo
optimizado.

\begin{table}[htbp]
\centering
\caption{Comparativa: SVM base vs.\ SVM optimizada
  ($\tau_{\text{des}} = 75\%$, test = E3).}
\label{tab:refinamiento}
\footnotesize
\begin{tabular}{@{}lrrrr@{}}
\toprule
Modelo & {$F_1$} & {Acc.} & {Adj.} & {MAE} \\
\midrule
SVM base (RBF, $t = 0{,}50$)   & 0{,}36 & 0{,}503 & 0{,}859 & 0{,}638 \\
SVM opt.\ (lineal, $t = 0{,}30/0{,}25$) & 0{,}40 & 0{,}602 & 0{,}985 & --- \\
\midrule
$\Delta$ & {$+0{,}04$} & {$+0{,}099$} & {$+0{,}126$} & --- \\
\bottomrule
\end{tabular}
\end{table}

La configuración óptima (SVM lineal, $C = 5$, $k = 10$ \textit{features}
por información mutua, pesos de clase $\{0{:}5,\; 1{:}1,\; 2{:}5\}$,
umbrales de decisión $t_1 = 0{,}30$, $t_2 = 0{,}25$) produce una mejora
de $+12{,}6$ puntos porcentuales en exactitud adyacente (de 85{,}9\% a
98{,}5\%) y $+9{,}9$ puntos en exactitud exacta (de 50{,}3\% a 60{,}2\%).
La matriz de confusión (Tabla~\ref{tab:cm_opt}) muestra que el modelo
optimizado logra discriminar parcialmente las tres clases, aunque la
clase ``sin desgaste'' sigue siendo la más difícil (5{,}3\% de
clasificación correcta), ya que la mayoría de sus muestras se clasifican
como ``medianamente desgastado'' (error de un paso ordinal).

\begin{table}[htbp]
\centering
\caption{Matriz de confusión del modelo SVM optimizado
  ($\tau_{\text{des}} = 75\%$, test = E3).}
\label{tab:cm_opt}
\footnotesize
\begin{tabular}{@{}l*{3}{r}@{}}
\toprule
& \multicolumn{3}{c}{Predicción} \\
\cmidrule(lr){2-4}
Real & {Sin desg.} & {Med.\ desg.} & {Desg.} \\
\midrule
Sin desgaste ($n = 95$) & 5 & 81 & 9 \\
Med.\ desgastado ($n = 344$) & 0 & 302 & 42 \\
Desgastado ($n = 144$) & 0 & 100 & 44 \\
\bottomrule
\end{tabular}
\end{table}

Tres factores explican la mejora: (1)~los pesos de clase asimétricos
$\{0{:}5,\; 1{:}1,\; 2{:}5\}$ fuerzan al clasificador a discriminar las
clases minoritarias en lugar de predecir la clase mayoritaria por defecto;
(2)~la calibración de los umbrales de decisión ($t_1 = 0{,}30$ vs.\
$0{,}50$) corrige la tendencia del modelo base a subestimar las clases
extremas; y (3)~las 7 \textit{features} derivadas (ratios espectrales,
indicadores de estabilidad tonal, energía logarítmica) proporcionan
representaciones más linealmente separables, favoreciendo al kernel lineal.

% ── 3.8 Generalización cruzada ─────────────────────────────────────────────
\subsection{Generalización cruzada entre experimentos}

La Fig.~\ref{fig:cross_exp} presenta el heatmap de generalización
\textit{leave-one-experiment-out} con umbral al 97\%.

\begin{figure}[htbp]
  \centering
  \includegraphics[width=\columnwidth]{cross_experiment_heatmap.pdf}
  \caption{Generalización \textit{leave-one-experiment-out}
    (umbral 97\%). Cada fila es un experimento evaluado como test.}
  \label{fig:cross_exp}
\end{figure}

Los resultados revelan tres grupos:
\begin{itemize}
  \item \textbf{Alta generalización} (E5, E6): $F_1 > 0{,}63$, exactitud
    $> 83\%$. Son ensayos progresivos de solo dos clases (sin y
    medianamente desgastado), más fáciles de clasificar.
  \item \textbf{Generalización moderada} (E1, E3, E4, E7):
    $F_1 \in [0{,}31;\; 0{,}44]$. Representan condiciones típicas con
    tres clases.
  \item \textbf{Falla de generalización} (E2): $F_1 = 0{,}18$,
    exactitud $= 22\%$. E2 usa micrófono dinámico mientras el resto del
    entrenamiento es condensador, generando un \textit{domain shift}
    severo.
\end{itemize}

Un hallazgo positivo es que la \textbf{exactitud adyacente es
$\geq 97\%$ en todos los experimentos}, incluyendo E2. Esto significa
que incluso cuando el modelo falla, sus errores son de un solo paso
ordinal.

% ============================================================================
\section{Discusión}
\label{sec:discusion}

% ── 4.1 Interpretación física ───────────────────────────────────────────────
\subsection{Interpretación física de la sensibilidad al umbral}

La curva de sensibilidad al umbral (Fig.~\ref{fig:threshold}) tiene una
interpretación física directa. El desgaste severo (cercano al fallo)
produce fenómenos acústicos identificables: \textit{chatter} regenerativo,
\textit{chirping} por fricción en el flanco desgastado y aumento de la
emisión acústica por deformación plástica severa~\cite{bhuiyan2016}. Estos
fenómenos se manifiestan como cambios abruptos en el contraste espectral,
la tasa de \textit{onsets} y el factor de cresta --- precisamente las
\textit{features} más discriminativas según el estudio de ablación.

En contraste, el desgaste temprano-intermedio (60--75\% de vida útil) se
caracteriza por desgaste abrasivo gradual del flanco que modifica las
fuerzas de corte y la rugosidad superficial, pero genera cambios acústicos
sutiles que se confunden con la variabilidad normal del proceso. Esto
explica por qué el $F_1$ crece monótonamente con el umbral: las
\textit{features} \textit{hand-crafted} son descriptores estadísticos
globales del espectro que capturan cambios bruscos pero no transiciones
graduales. La caída pronunciada del $F_1$ al 60\% ($F_1 = 0{,}30$ en SVM)
se debe a que, con un umbral tan bajo, la clase ``desgastado'' incluye
muestras que corresponden a desgaste moderado cuya firma acústica es
prácticamente indistinguible de la clase ``medianamente desgastado'',
generando una frontera de decisión ruidosa que incrementa tanto los falsos
positivos como los falsos negativos de la clase minoritaria.

% ── 4.2 Relevancia operacional ──────────────────────────────────────────────
\subsection{Relevancia operacional de la exactitud adyacente}

La exactitud adyacente es particularmente relevante para el monitoreo
industrial porque cuantifica la severidad de los errores. Un sistema que
clasifica una broca ``medianamente desgastada'' como ``sin desgaste''
comete un error de un paso: el operario no cambiará la herramienta pero
aún tiene margen antes del fallo. En cambio, clasificar ``sin desgaste''
como ``desgastado'' (error de dos pasos) provocaría un cambio innecesario
de herramienta con impacto en productividad.

Los resultados muestran que el error de dos pasos es \textit{infrecuente}:
a umbrales $\geq 85\%$, la exactitud adyacente alcanza $\geq 92{,}6\%$
(SVM) y $\geq 96{,}9\%$ (RF), y es raro incluso a umbrales bajos. Esto
sugiere que un sistema de monitoreo basado en estas \textit{features},
aunque no alcance alta precisión por clase, es \textit{seguro} en el
sentido de que sus errores no son catastróficos.

% ── 4.3 Micrófono como covariable ──────────────────────────────────────────
\subsection{El tipo de micrófono como fuente de \textit{domain shift}}

El resultado más preocupante de la generalización cruzada es la caída
dramática en E2 ($F_1 = 0{,}18$), el único ensayo con micrófono dinámico
en el set de evaluación mientras el entrenamiento se realiza
predominantemente con micrófono condensador. Dado que los micrófonos
dinámicos tienen menor sensibilidad y rango de frecuencia
(50~Hz--15~kHz vs.\ 40~Hz--20~kHz), las \textit{features} espectrales
cambian sistemáticamente aunque el proceso de corte sea el mismo.

Este hallazgo tiene implicaciones prácticas: un sistema de monitoreo
acústico debe calibrarse o adaptarse al tipo de transductor utilizado.
Las técnicas de adaptación de dominio~\cite{banda2024} o la inclusión
del tipo de micrófono como covariable podrían mitigar este efecto.

Los resultados del preprocesamiento con \textit{spectral gating}
(Sección~\ref{sec:noisereduce}) aportan evidencia directa sobre esta
hipótesis. Al aplicar reducción de ruido, el $F_1$ macro en E2 pasa de
0{,}18 a 0{,}33 ($+83{,}9\%$ relativo) y la exactitud de 23{,}4\% a
58{,}6\%, sin degradar el rendimiento en test (E3, mismo tipo de
micrófono que el entrenamiento). Esto sugiere que una fracción
sustancial del \textit{domain shift} entre tipos de micrófono se
origina en diferencias del piso de ruido captado por cada transductor
---componente que el \textit{spectral gating} atenúa al normalizar el
contenido no informativo de la señal. El efecto residual ($F_1 = 0{,}33$
aún inferior al test $F_1 = 0{,}34$) indica que persisten diferencias
intrínsecas de respuesta en frecuencia que requieren estrategias
complementarias (adaptación de dominio, calibración por transductor o
inclusión del tipo de micrófono como covariable en el modelo).

% ── 4.4 Comparación con la literatura ──────────────────────────────────────
\subsection{Comparación con la literatura}

La Tabla~\ref{tab:literatura} contextualiza los resultados frente a
trabajos previos. Las comparaciones son aproximadas porque difieren en
materiales, sensores y definición de desgaste; sin embargo, las filas de
Orejarena~(2014) y este trabajo comparten el mismo conjunto de datos de
taladrado, lo que permite una comparación directa de metodologías.

\begin{table*}[htbp]
\centering
\caption{Comparación con trabajos previos en monitoreo acústico de
  herramientas de corte. Filas con $\dagger$ comparten el mismo
  dataset experimental (Orejarena, 2014).}
\label{tab:literatura}
\footnotesize
\begin{tabular}{@{}p{2.8cm}p{1.6cm}p{2.0cm}p{0.8cm}p{1.0cm}p{1.3cm}p{1.6cm}@{}}
\toprule
Referencia & Sensor & Método & Pre-proc. & Clases & Exactitud & Adj.\ acc. \\
\midrule
Banda et al.\ (2024)     & Acelerómetro  & SVM + SVD         & No  & 3         & 97\%      & N/D \\
Bhuiyan et al.\ (2016)   & AE + mic      & ANN/SVM           & No  & 2         & 94\%      & N/D \\
Charoenprasit (2018)     & Mic husillo   & SVM               & No  & 2         & 92\%      & N/D \\
Orejarena (2014)$^\dagger$ & 1 mic (din.) & Corr.\ + SVM    & No  & 2         & 85\%      & N/D \\
\midrule
\textbf{Este trabajo (base)}$^\dagger$ & \textbf{2 mic}   & \textbf{SVM/RF ordinal}  & \textbf{Sí}  & \textbf{3 (ord.)} & \textbf{58--71\%$^a$} & \textbf{98{,}6--99{,}5\%} \\
\textbf{Este trabajo (opt.)}$^\dagger$  & \textbf{2 mic}   & \textbf{SVM lineal opt.} & \textbf{Sí}  & \textbf{3 (ord.)} & \textbf{60\%$^b$}     & \textbf{98{,}5\%} \\
\bottomrule
\multicolumn{7}{@{}l}{\footnotesize $^a$Exactitud exacta al 97\% (SVM: 58\%, RF: 71\%). N/D = no disponible (clasificación binaria, adj.\ acc.\ $\equiv$ exactitud).} \\
\multicolumn{7}{@{}l}{\footnotesize $^b$SVM lineal, $C=5$, $k=10$ \textit{features}, umbrales $t_1{=}0{,}30$/$t_2{=}0{,}25$, umbral de etiquetado 75\%.}
\end{tabular}
\end{table*}

La exactitud exacta de este trabajo es inferior a la de estudios previos,
pero la comparación es inequitativa por tres razones: (1)~se resuelve un
problema de \textit{tres clases ordinales}, no binario --- una tarea
inherentemente más difícil donde Orejarena~(2014), con el mismo dataset
pero formulación binaria, alcanza solo 85\%; (2)~se evalúa con
separación estricta por experimento (\textit{GroupKFold}), no con
partición aleatoria que infla resultados; (3)~la métrica relevante para
problemas ordinales es la \textbf{exactitud adyacente}: con 98{,}6\% en
la configuración base y 98{,}5\% en la optimizada, este trabajo es
competitivo con todos los enfoques de la literatura, incluidos los
binarios, demostrando que los errores del sistema nunca son
catastróficos desde el punto de vista operacional.

% ── 4.5 Limitaciones ───────────────────────────────────────────────────────
\subsection{Limitaciones}

Este estudio presenta varias limitaciones que deben considerarse al
interpretar los resultados:

\begin{enumerate}
  \item \textbf{Tamaño del conjunto ``desgastado''}: solo 16 muestras
    genuinas de fallo catastrófico (sin augmentación), lo que introduce
    variabilidad alta en las métricas (IC~95\% de $F_1$:
    $[0{,}38;\; 0{,}75]$ a 97\%).
  \item \textbf{Ausencia de medición de VB}: las etiquetas se basan en
    porcentaje de vida útil, no en desgaste de flanco medido según
    ISO~8688. Esto introduce incertidumbre en la frontera entre clases.
  \item \textbf{\textit{Domain shift} por micrófono}: el modelo no
    generaliza al micrófono dinámico (E2), limitando su aplicabilidad a
    configuraciones con micrófono condensador.
  \item \textbf{Un solo material y diámetro por test}: el test set (E3)
    solo contiene brocas de 6~mm en AISI~4140.
  \item \textbf{Patrón G-code como confusión potencial}: aunque no se
    considera variable explicativa, E3 (lineal) difiere de E4 (espiral).
\end{enumerate}

% ============================================================================
\section{Conclusiones}
\label{sec:conclusiones}

Se presentó un sistema de clasificación ordinal del desgaste de brocas CNC
mediante análisis acústico, con un análisis sistemático de sensibilidad al
umbral de etiquetado. Las principales conclusiones son:

\begin{enumerate}
  \item Las \textit{features} acústicas \textit{hand-crafted} detectan
    eficazmente el fallo catastrófico ($\geq 97\%$ vida útil:
    $F_1 = 0{,}46$, adj.\ acc.\ $= 98{,}6\%$) pero tienen capacidad
    limitada para el desgaste progresivo temprano ($\geq 60\%$:
    $F_1 = 0{,}30$, adj.\ acc.\ $= 83{,}4\%$).

  \item La exactitud adyacente permanece $\geq 83\%$ en todo el rango de
    umbrales, garantizando que los errores nunca exceden un paso ordinal.
    Esto hace al sistema \textit{seguro} para monitoreo industrial.

  \item El umbral de 85\% representa el mejor compromiso operacional:
    RF logra adj.\ acc.\ $= 96{,}9\%$ con 86 muestras de test (vs.\ 18
    a 97\%), proporcionando soporte estadístico más robusto.

  \item Un subconjunto de 7 \textit{features} supera al conjunto completo
    de 10 en todos los umbrales ($F_1$: $+0{,}01$ a $+0{,}04$), indicando
    que tres \textit{features} (\texttt{onset\_rate},
    \texttt{spectral\_flatness}, \texttt{duration\_s}) introducen ruido.

  \item La etapa de refinamiento (ingeniería de \textit{features}
    derivadas, calibración de hiperparámetros y umbrales de decisión)
    eleva la exactitud adyacente de la SVM al 75\% de 85{,}9\% a 98{,}5\%
    y la exactitud exacta de 50{,}3\% a 60{,}2\%, demostrando que la
    calibración del clasificador ordinal es tan importante como la
    selección del modelo.

  \item El tipo de micrófono es la principal fuente de \textit{domain
    shift}, con degradación severa al evaluar con micrófono dinámico
    ($F_1 = 0{,}18$) vs.\ condensador ($F_1 = 0{,}33$).

  \item La reducción de ruido espectral (\textit{spectral gating})
    mejora la exactitud adyacente en test de 84{,}7\% a 90{,}1\% y
    reduce drásticamente el \textit{domain shift} entre tipos de
    micrófono ($F_1$: 0{,}18 $\rightarrow$ 0{,}33; exactitud: 23\%
    $\rightarrow$ 59\% en validación), demostrando que una fracción
    significativa de la variabilidad inter-transductor se origina en
    diferencias del piso de ruido.
\end{enumerate}

El bajo costo computacional del enfoque propuesto (inferencia $< 1$~ms
por muestra con 10 \textit{features} y SVM lineal) lo hace viable para
sistemas de monitoreo en línea embebidos, sin requerir GPU ni
infraestructura de cómputo especializada. Esta característica es relevante
para la adopción en entornos de manufactura donde los recursos
computacionales son limitados.

Trabajo futuro incluye: (1)~explorar representaciones espectrales
aprendidas (CNN sobre mel-espectrogramas) para capturar transiciones
graduales de desgaste; (2)~incorporar medición de VB con microscopio
calibrado según ISO~8688; (3)~profundizar en técnicas de adaptación de
dominio complementarias al \textit{spectral gating} para cerrar la brecha
residual entre tipos de micrófono; (4)~validar el sistema con brocas de
diferente geometría (Dormer~A114) y condiciones de refrigeración
variables; (5)~evaluar la escalabilidad del enfoque a otros procesos de
corte (fresado, torneado) donde la emisión acústica presenta
características similares al taladrado.

% ============================================================================
% AGRADECIMIENTOS
\section*{Agradecimientos}

Los autores agradecen a la Escuela de Ingeniería Mecánica de la
Universidad Industrial de Santander por facilitar el acceso al Laboratorio
de Mecanizado y al centro de mecanizado CNC Leadwell V-20.

% ============================================================================
% REFERENCIAS
\begin{thebibliography}{99}

\bibitem{groover2007}
M.~P.~Groover, \textit{Fundamentos de manufactura moderna: materiales,
procesos y sistemas}, 3a ed. México: McGraw-Hill, 2007.

\bibitem{bhuiyan2016}
M.~S.~H.~Bhuiyan, I.~A.~Choudhury y M.~Dahari, ``Application of acoustic
emission sensor to investigate the frequency of tool wear and
\textit{tool failure} in turning operation,'' \textit{Measurement},
vol.~92, pp.~208--217, 2016.

\bibitem{charoenprasit2018}
K.~Charoenprasit, P.~Srikummoon y S.~Sombatsiri, ``Monitoring tool wear
in CNC milling using spindle motor noise,'' \textit{Int.\ J.\ Mech.\
Eng.\ Robot.\ Res.}, vol.~7, no.~5, pp.~522--527, 2018.

\bibitem{orejarena2014}
N.~Orejarena Osorio, ``Estudio de la señal acústica en función de la
vida de la herramienta en un proceso de taladrado,'' Tesis M.Sc.,
Universidad Industrial de Santander, 2014.

\bibitem{banda2024}
T.~Banda, A.~Ali, J.~Asku y H.-C.~Möhring, ``Multi-sectional SVD-based
machine learning for tool wear assessment and prediction in turning
process,'' \textit{Int.\ J.\ Adv.\ Manuf.\ Technol.}, vol.~132,
no.~7--8, pp.~3585--3600, 2024.

\bibitem{iso8688}
ISO~8688-1:1989, \textit{Tool life testing in milling --- Part 1: Face
milling}, International Organization for Standardization, 1989.

\bibitem{frank2001}
E.~Frank y M.~Hall, ``A simple approach to ordinal classification,'' in
\textit{Proc.\ European Conf.\ Machine Learning (ECML)}, Springer,
2001, pp.~145--156.

\bibitem{macphail2024}
V.~MacPhail, J.~Buxton, F.~Bussière y A.~C.~Smith, ``Audio compression
affects acoustic indices and classifications of ecological sounds,''
\textit{Bioacoustics}, vol.~33, no.~1, pp.~1--22, 2024.

\bibitem{montgomery2013}
D.~C.~Montgomery, \textit{Design and Analysis of Experiments}, 8th~ed.
Hoboken, NJ: Wiley, 2013.

\bibitem{meneses2020}
J.~E.~Meneses Flórez y N.~Peña, ``Estudio correlacional entre la señal
acústica y el desgaste de la herramienta en un proceso de taladrado,''
\textit{Rev.\ Colomb.\ Tecnol.\ Avanzada (RCTA)}, vol.~2, no.~34,
pp.~1--8, 2020.

\bibitem{geron2023}
A.~Géron, \textit{Hands-on Machine Learning with Scikit-Learn, Keras,
and TensorFlow}, 3rd~ed. Sebastopol, CA: O'Reilly, 2023.

\bibitem{huang2024}
X.~Huang, W.~Wang y S.~Li, ``Time series feature selection method based
on mutual information,'' \textit{Applied Sciences}, vol.~14, no.~5,
art.~1960, 2024.

\bibitem{li2024}
Y.~Li, J.~Li, Y.~Zhang y Z.~Zhang, ``The Shapley value: from
cooperative game theory to explainable AI,'' \textit{Autonomous
Intelligent Systems}, vol.~4, art.~2, 2024.

\bibitem{sainburg2020}
T.~Sainburg, M.~Thielk y T.~Q.~Gentner, ``Finding, visualizing, and
quantifying latent structure across diverse animal vocal repertoires,''
\textit{PLOS Computational Biology}, vol.~16, no.~10, art.~e1008228,
2020. [Biblioteca \texttt{noisereduce}: \url{https://github.com/timsainb/noisereduce}]

\end{thebibliography}

\end{document}










